% !TeX root = ../main.tex
\section{The Bestvina-Brady Theorem for Flag Complexes} 
\label{sec:L06-W03D1}

\begin{definition}[flag simplicial complex]
    \label{def:flag-simplicial-complex}
    A simplicial complex \(\G\) is \emph{flag} if whenever a set of vertices \(\{v_0, \dots, v_k\}\) span a complete subgraph, then they span a \(k\)-simplex.
\end{definition}

In other words, every non-simplex contains a non-edge.

Any graph \(\G\) determines a flag simplicial complex.

\begin{definition}[RAAGs of flag simplicial complex]
    \label{def:flag-RAAG}
    Let \(\G\) be a simplicial graph with vertex set \(V\) and edge set \(E\). Let \(A_\G\) be the associated RAAG (See \cref{def:RAAG}). If \(L\) is a flag simplicial complex, Then \(A_L:= A_{L^{(1)}}\).

    Define \(\phi_L:A_L \to \Z\) by \(v\mapsto 1\). Let \(K_L = \ker \phi_L\).
\end{definition}

\begin{remark}
    Since \(\phi_L\) is a homomorphism, \(K_L\) is related to the ``topology'' of flag of  \(\G\).
\end{remark}

\begin{example}
    \label{ex:flag-RAAG}
    Let \(L\) be a square, then \(A_L = \mathbb{F}_2 \times \mathbb{F}_2\).

    In general, join of graphs corresponds to product of groups.(Disjoint union corresponds to free product).
\end{example}


\begin{definition}
    \label{def:cyclic-space}
    A space \(X\) is \emph{acyclic} if \(\tilde H_i (X) = 0\) for all \(i \geq 0\).
\end{definition}


Recall that a group is \emph{perfect} if it has trivial abelianization, i.e. \(\tilde H_1(X) \) is trivial. By definition, every acyclic group is perfect.

\begin{definition}
    \label{def:homologically-k-connected}
    A space \(X\) is \emph{homologically \((k)\)-connected} if \(\tilde H_i (X)=0\) for \(i \leq k\)
\end{definition}

\begin{restatable}[The Bestvina-Brady Theorem]{theorem}{BestvinaBrady}
    \label{thm:Bestvina-Brady-finiteness-of-flag-RAAG}
    Let \(L\) be a finite flag simplicial complex. Let \(A_L\) be the associated RAAG. Let \(\phi_L : A_L \to \Z\) be as in \cref{def:flag-RAAG} and let \(K_L = \ker \phi_L\).
    \begin{enumerate}
        \item \(K_L\) is finitely presented if and only if \(L\) is simply connected.
        \item \(K_L\) has type \(FP\) if and only if \(L\) is acyclic.
        \item \(K_L\) has type \(FP_n\) if and only if \(L\) is homologically \((n-1)\)-connected.
    \end{enumerate} 
\end{restatable}


\begin{theorem}[{\cite[Theorem 4.32]{hatcher_algebraic_2001}}]
    \label{thm:Hurwicz-homologically-connected-and-homotopy-group}
    If \(X\) is \((n-1)\)-connected, \(n\geq 2\), then \(\tilde H_i (X)= 0\) for \(0\leq i \leq n-1\), and \(H_n(X) \cong \pi_n(X)\).
\end{theorem}



\cref{thm:Hurwicz-homologically-connected-and-homotopy-group} will guide our search for a non-contractible acyclic space. We also want \(\pi_1 \) to be perfect.

\begin{example}[Poincare's homology sphere]
    \label{ex:Poincare-homology-sphere}
    Take an empty dodecahedron such that opposite sides are the same up to a twist. Identify opposite sides with a \(\pi/5\) twist, then we get the \(2\)-skeleton of Poincare's homology sphere, denoted as \(X\). 

    
    Let \(M\) be the homology sphere built from a solid dodecahedron (which is a 3-manifold, and has \(X\) as its \(2\)-skeleton). We have  \(\tilde H_3(M) \cong \tilde H_3(\mathbb{S}^3)= \Z\), and that \(\tilde H_i(M) \cong \tilde H_i(\mathbb{S}^3)= 0\)  for all other \(i\)'s.
   
    Recall Universal coefficient: if \(H_i(M)\) is torsion free, then \(H^{i+1}(M;\Z ) \cong \Hom{H_i(M)}{\Z}{}\).

    Recall     Poincare Duality: \(H_i(M) \cong H^{3-i}(M;\Z)\)

    Since \(H_0(M) = 0\), \(H_1(M)=0 \implies H^1(M;\Z) = 0\) by Universal Coefficient theorem, then \(H_2(M)=0\) (by Poincare Duality).

    So do Cellular chain complex for \(M\) we get 
    \[0 \to C_3(M) \xrightarrow{\a_3} C_2(M) \to C_1(M)\to \Z \to 0\]

    Observe that \(\a_3 =0\) because \(M\) has  no boundary after gluing. Also \(C_3(M)  =\Z\) because there's only one \(3\)-cell.

    We can ignore \(C_3(M)\). The rest of the chain (and the homology) coincides with that of \(X\): 
    \[ 0 \to C_2(X) \to C_1(X) \to \Z \to 0\]

    And so all homologies of \(X\) are trivial. 

    This \(X\) with \(K_X\) gives the first example of a group which is acyclic but non-contractible complex.
    
    Now subdivide each face to (10) triangles to obtain a flag complex \(L\). This \(L\) is not simply connected, so \(K_L\) is not finitely presented. It is still acyclic, so \(K_L\) has type \(FP\). This gives the first example of a group that has \(FP\) but is not finitely presented.
\end{example}

\begin{remark}[Seifert-Weber space]
    \label{rmk:Seifert-Weber-space}
    If we do \(3\pi/5\) rotation with filled dodecahedron, then we get an infinite perfect \(\pi_1\), and it's hyperbolic. The obtained space is known as Seifert-Weber space.

    One can see that by seeing how there are  \(5\) pentagons around an edge. An angle argument suggests the hyperbolic geometry.
\end{remark}


Our goal now is to prove \cref{thm:Bestvina-Brady-finiteness-of-flag-RAAG}. We first introduce yet another finiteness property.

\begin{restatable}{definition}{TypeFH}
    \label{def:type-FHn-finiteness-property}
    A group \(G\) has type \(FH_n\) (resp. \(FH\)) if there exists a homologically \((n-1)\)-connected (resp. acyclic) complex \(X\) with a free cellular cocompact \(G\) action.
\end{restatable}


\begin{lemma}
    \label{lem:FPn-implies-FFn}
    \(FP_n \implies FF_n\).
\end{lemma}
\begin{proof}
    Let \(M\) be a projective \(\Z G\)-module.
    Take \(M\) as a quotient of free module:
    \[ 0 \to K \to F \to M \to 0 . \]
    If \(M\) is finitely generated, we can assume \(F\) is also finitely generated. By \cref{prop:characterization-of-projctive-module} \(K\) is projective and the sequence splits, which forces \(K\) to be finitely generated. Hence \(F=K\oplus M\).

    Take a finitely generated projective resolution \(\{P_i\}_{i=0}^n\)  of \(\Z\). Then there exists \(F_0 =P_0\oplus Q_0\). Build:
    \[ P_1\oplus Q_0 \to P_0 \oplus Q_0 \to \Z \]

    As a direct sum of projective modules, \(P_1 \oplus Q_0\) is also projective. We can then find \(Q_1\) such that \(Q_1 \oplus P_1\oplus Q_0 \) is a finitely generated free module. We can repeat this procedure to eventually obtain a finitely generated free resolution.
\end{proof}


We now have four types of finiteness properties: \(F, FP, FF, FH \) (See \cref{def:finite-type-KG1,def:type-FF-FP,def:type-FHn-finiteness-property}).  For a collection of relationships between these properties, see \cref{sec:collection-of-types}.

